\section{实验方法}

\subsection{sgRNA 设计与病毒载体构建}
\subsubsection{LentiV2-puro 载体限制性酶切反应}
\begin{enumerate}
    \item 在 \qty{0.2}{\micro\liter} 离心管中配制反应体系：
    
    \begin{table}[H]
        \centering
        \begin{tabular}{C{4cm} C{4cm}}
            \toprule
            \textbf{试剂} & \textbf{体系} \\
            \midrule
            LentiV2-puro        & \qty{5}{\micro\gram} \\
            Bsmb1               & \qty{5}{\micro\liter} \\
            CutSmart\textregistered\ buffer & \qty{5}{\micro\liter} \\
            DTT                 & \qty{1}{\micro\liter} \\
            ddH$_2$O            & X\,\si{\micro\liter} \\
            \midrule
            \textbf{Total}      & \textbf{\qty{50}{\micro\liter}} \\
            \bottomrule
        \end{tabular}
    \end{table}
    
    \item 混匀 PCR 管，然后短暂离心；将 PCR 管置于 \qty{37}{\celsius} 水浴或者金属浴中 \qty{4}{\hour}。
    \item 1\% 琼脂糖凝胶电泳鉴定。
\end{enumerate}

\subsubsection{DNA 凝胶电泳}

\begin{enumerate}
    \item 取洁净有机玻璃制胶槽，水平放置并安装挡板，确保装置密封防止渗漏。将样品梳插入制胶槽相应卡槽中备用。
    \item 称取 \qty{1.0}{\g} 琼脂糖粉末置于 \qty{250}{\ml} 锥形瓶中，加入 \qty{100}{\ml} 1\(\times\)TAE 电泳缓冲液，轻轻摇匀后使用微波炉或水浴加热至琼脂糖完全溶解，溶液应清澈透明无颗粒。待溶液不冒泡后加入 \qty{10}{\ul} GoldView 核酸染料。将溶解的琼脂糖溶液冷却至约 \qtyrange{60}{65}{\celsius}（手感温热但不烫手），沿玻璃棒缓缓倒入制胶槽，避免产生气泡。室温静置约 30 分钟使凝胶完全凝固，待凝固后小心拔出样品梳，可见清晰加样孔形成。
    \item 将凝固的凝胶连同制胶槽一并放入电泳槽中，加入 1\(\times\)TAE 电泳缓冲液至液面完全覆盖凝胶表面约 \qtyrange{1}{2}{\mm}。使用移液器将混合液缓慢加入加样孔中，注意勿刺破凝胶孔壁或溢出孔外。
    \item 盖上电泳槽盖，正确连接电极。接通电源，设定恒定电压 \qty{180}{\V} 进行电泳。当溴酚蓝指示剂迁移至凝胶全长 2/3\,--\,3/4 处时停止电泳（约 20\,--\,30 分钟）。
    \item 跑胶结束后关闭电源，取出凝胶，置于紫外凝胶成像仪或紫外透射仪中观察。在紫外光激发下，DNA 条带呈现明亮荧光，根据需要拍摄记录电泳结果。将目的片段小心切下，置于干净的 \qty{1.5}{\ml} EP 管中，进行凝胶回收。
\end{enumerate}

